5. Encuentre los siguientes conjuntos. \\
a) \( \bigcap_{n \in \mathbb{N}}\left[-\frac{1}{n}, 1+\frac{1}{n}\right] \). \\
Sean los intervalos
\[
I_n = \left[-\frac{1}{n}, 1 + \frac{1}{n}\right], \quad n \in \mathbb{N}.
\]

Queremos mostrar que $I_n \supseteq I_{n+1} \text{ para todo } n \in \mathbb{N}$, esto para usar propiedades de conjuntos.

\medskip

Observemos los extremos de cada intervalo:

\[
-\frac{1}{n} \leq -\frac{1}{n+1}, \quad \text{y} \quad 1 + \frac{1}{n} \geq 1 + \frac{1}{n+1}.
\]

Esto se cumple porque:

\[
\frac{1}{n} > \frac{1}{n+1} > 0,
\]

por lo que

\[
-\frac{1}{n} < -\frac{1}{n+1} \quad \Rightarrow \quad -\frac{1}{n} \leq -\frac{1}{n+1},
\]

y

\[
1 + \frac{1}{n} > 1 + \frac{1}{n+1}.
\]

Por lo tanto,

\[
I_n = \left[-\frac{1}{n}, 1 + \frac{1}{n}\right] \supseteq \left[-\frac{1}{n+1}, 1 + \frac{1}{n+1}\right] = I_{n+1}.
\]

\medskip

Como los intervalos forman una cadena decreciente:

\[
I_1 \supseteq I_2 \supseteq I_3 \supseteq \cdots,
\]

la intersección

\[
\bigcap_{n=1}^\infty I_n = \lim_{n \to \infty} I_n = \left[\lim_{n \to \infty} -\frac{1}{n}, \lim_{n \to \infty} \left(1 + \frac{1}{n}\right)\right] = [0,1].
\]


b) $\bigcap_{n \in \mathbb{N}} \left(1 - \frac{1}{n}, 1 + \frac{1}{n}\right)$

\medskip

Sea la sucesión de intervalos abiertos
\[
J_n = \left(1 - \frac{1}{n}, 1 + \frac{1}{n}\right), \quad n \in \mathbb{N}.
\]

Mostremos que J_n \supseteq J_{n+1} \text{ para todo } n \in \mathbb{N}.

\medskip

Observemos los extremos:

\[
1 - \frac{1}{n} < 1 - \frac{1}{n+1}, \quad \text{y} \quad 1 + \frac{1}{n} > 1 + \frac{1}{n+1}.
\]

Esto se debe a que \( \frac{1}{n} > \frac{1}{n+1} > 0 \), por lo que:

\[
1 - \frac{1}{n} < 1 - \frac{1}{n+1} \quad \Rightarrow \quad \text{el extremo izquierdo de } J_n \text{ es menor que el de } J_{n+1},
\]
y
\[
1 + \frac{1}{n} > 1 + \frac{1}{n+1} \quad \Rightarrow \quad \text{el extremo derecho de } J_n \text{ es mayor que el de } J_{n+1}.
\]

Por lo tanto,

\[
J_n \supseteq J_{n+1}.
\]

\medskip

Como los intervalos forman una cadena decreciente, la intersección es:

\[
\bigcap_{n=1}^\infty J_n = \lim_{n \to \infty} J_n = \left(\lim_{n \to \infty} \left(1 - \frac{1}{n}\right), \lim_{n \to \infty} \left(1 + \frac{1}{n}\right) \right) = (1, 1).
\]

\medskip

\text{Pero } (1,1) \text{ es el intervalo vacío, entonces:}

\[
\bigcap_{n=1}^\infty \left(1 - \frac{1}{n}, 1 + \frac{1}{n}\right) = \varnothing.
\]

\bigskip

c) \displaystyle \bigcup_{n \in \mathbb{N}} [2n, 2n + 1]

\medskip

Consideremos la sucesión de intervalos cerrados

\[
K_n = [2n, 2n + 1], \quad n \in \mathbb{N}.
\]

Cada intervalo tiene longitud 1 y comienza en un número par creciente.

\medskip

Observemos que para \( m \neq n \), los intervalos \( K_n \) y \( K_m \) son disjuntos porque

\[
\text{si } m > n, \quad \min K_m = 2m \geq 2(n+1) > 2n + 1 = \max K_n.
\]

\medskip

Por lo tanto, la unión

\[
\bigcup_{n=1}^\infty [2n, 2n + 1]
\]

es la unión disjunta de segmentos separados que son:

\[
[2,3] \cup [4,5] \cup [6,7] \cup \cdots.
\]

\medskip

No se puede simplificar más, y el conjunto resultante es

\[
\{ x \in \mathbb{R} : x \in [2n, 2n+1] \text{ para algún } n \in \mathbb{N} \}.
\]